\chapter{Material Module Architecture Analysis}
\label{ch:material-module-analysis}

\section{Purpose}

This chapter provides a focused architectural analysis of the current
\texttt{materials} module.  The goal is not merely to rename classes or to
standardise terminology, but to identify the semantic boundaries that must be
made explicit if \texttt{fall\_n} is to evolve toward:
\begin{itemize}
  \item reinforced concrete modelling at several scales;
  \item materially nonlinear and dynamic analysis;
  \item finite-kinematics structural and continuum formulations;
  \item continuum-scale cracking, damage, localisation, and fracture.
\end{itemize}

The analysis is guided by three constraints already visible across the code
base:
\begin{itemize}
  \item the local element kernels must remain inexpensive and friendly to
        compile-time optimisation;
  \item the library must preserve the current genericity across continuum,
        beam, and shell formulations;
  \item the constitutive layer must remain reusable beyond post-processing,
        because the same reconstruction policies will later be needed for
        scale transfer and multiscale initialisation.
\end{itemize}

\section{Current Architectural Layers}

\subsection{Measure spaces}

At the lowest level, the module already contains a useful abstraction:
\begin{itemize}
  \item \texttt{Strain<N>} and \texttt{Stress<N>} for continuum Voigt spaces;
  \item \texttt{BeamGeneralizedStrain<N,Dim>} and \texttt{BeamSectionForces<N>};
  \item \texttt{ShellGeneralizedStrain<N,Dim>} and \texttt{ShellResultants<N>}.
\end{itemize}

The concepts in \texttt{ConstitutiveRelation.hh} correctly recognise that a
constitutive relation should fundamentally operate on a pair
\[
(\text{kinematic measure},\ \text{tensional conjugate}),
\]
independently of whether that pair belongs to a solid, a beam section, or a
shell section.  This is one of the strongest parts of the current design and
should be preserved.

\subsection{Relation classes}

The code currently places under the same umbrella:
\begin{itemize}
  \item continuum elastic relations such as
        \texttt{ContinuumIsotropicRelation};
  \item reduced section relations such as
        \texttt{TimoshenkoBeamSection3D} and \texttt{MindlinShellSection};
  \item inelastic local models such as
        \texttt{PlasticityRelation}, \texttt{KentParkConcrete}, and
        \texttt{MenegottoPintoSteel};
  \item reduction-based section models such as \texttt{FiberSection3D}.
\end{itemize}

All of them satisfy the same compile-time concept:
\[
\texttt{compute\_response}(k), \qquad
\texttt{tangent}(k),
\]
with optional inelastic extensions through
\[
\texttt{update}(k), \qquad
\texttt{internal\_state()}.
\]

This common contract is practical, but it already reveals a semantic
overloading: some of these types are truly \emph{material laws}, while others
are \emph{section laws} or even \emph{constitutive reductions}.

\subsection{Material instance and type erasure}

The next layer is \texttt{MaterialInstance<Relation,StatePolicy>}, which wraps:
\begin{itemize}
  \item a constitutive relation object;
  \item a per-point kinematic state container;
  \item access to the underlying relation through
        \texttt{constitutive\_relation()}.
\end{itemize}

This is then lifted to the type-erased \texttt{Material<Policy>} wrapper,
which stores:
\begin{itemize}
  \item the concrete \texttt{MaterialInstance};
  \item an update strategy (\texttt{ElasticUpdate} or
        \texttt{InelasticUpdate});
  \item export helpers such as \texttt{internal\_field\_snapshot()} and
        \texttt{section\_snapshot()}.
\end{itemize}

\subsection{Constitutive sites}

Finally, the constitutive object is attached to geometry through:
\begin{itemize}
  \item \texttt{MaterialPoint<Policy>} in continuum elements;
  \item \texttt{MaterialSection<Policy>} in beams and shells.
\end{itemize}

This layering already captures a valuable idea:
\begin{center}
\emph{geometry and topology are not constitutive objects;}\\
\emph{a constitutive object is bound to a sampling site.}
\end{center}

\section{What is Good and Should Be Preserved}

\subsection{The measure-based generic contract}

The most important success of the current module is that the constitutive
interface is not hard-wired to Cauchy stress or to infinitesimal continuum
strain.  Instead, it works on generic conjugate pairs.  This is exactly what a
library needs if it aims to support:
\begin{itemize}
  \item continuum solids;
  \item structural sections;
  \item future higher-order or enriched formulations;
  \item multiscale reductions and reconstructions.
\end{itemize}

\subsection{Static polymorphism at the relation level}

Concrete constitutive classes are template-friendly and concept-constrained.
This gives:
\begin{itemize}
  \item strong compile-time checking;
  \item direct inlining opportunities inside element kernels;
  \item no virtual dispatch when the full type is known.
\end{itemize}

This direction is correct and should remain the default in the hot path.

\subsection{Separation between geometry and constitutive sites}

The distinction between \texttt{IntegrationPoint} and
\texttt{MaterialPoint} / \texttt{MaterialSection} is also correct in spirit.
It is consistent with the earlier ontology work in the \texttt{Domain}
refactor:
\begin{itemize}
  \item the integration point is a numerical sampling site;
  \item the constitutive point is a physical role attached to that site.
\end{itemize}

\subsection{The reduction-based path for structural elements}

The existence of \texttt{FiberSection3D} is strategically important.  It proves
that the same element pipeline can host:
\begin{itemize}
  \item a linear reduced section law;
  \item a nonlinear reduced section law reconstructed from uniaxial
        constituents.
\end{itemize}

This is exactly the kind of reduction/reconstruction mechanism that will later
be needed for layered shells, smeared reinforcement, and scale-bridging
between frame and continuum descriptions.

\section{Observed Semantic Inconsistencies}

\subsection{The word \texttt{Material} currently means too many things}

Today the code uses the word \emph{material} for at least four different
concepts:
\begin{enumerate}
  \item a physical substance or phase
        (concrete, steel, confined concrete);
  \item a constitutive law
        (\texttt{KentParkConcrete}, \texttt{MenegottoPintoSteel});
  \item a reduced section law
        (\texttt{TimoshenkoBeamSection3D}, \texttt{FiberSection3D});
  \item a type-erased operational wrapper
        (\texttt{Material<Policy>}).
\end{enumerate}

This is the single most important semantic problem in the current module.
Architecturally, these should not be collapsed into one noun.

\subsection{\texttt{MaterialPolicy} is not a material}

\texttt{SolidMaterial<6>}, \texttt{BeamMaterial<6,3>}, and
\texttt{ShellMaterial<8,3>} do not describe a physical material.  They define
a \emph{constitutive space}:
\[
\mathcal{S} =
(\text{kinematic type},\ \text{conjugate type},\ \text{dimension}).
\]

So the name \texttt{MaterialPolicy} is semantically misleading.  The type is
really a \emph{measure-space policy} or a \emph{constitutive space descriptor}.
The current name propagates confusion upward into \texttt{Material},
\texttt{MaterialPoint}, and \texttt{MaterialSection}.

\subsection{\texttt{MaterialInstance} is not just a material instance}

\texttt{MaterialInstance<Relation,StatePolicy>} currently bundles:
\begin{itemize}
  \item a constitutive relation;
  \item kinematic state storage;
  \item direct access to the relation;
  \item optional internal-state access for inelastic models.
\end{itemize}

Semantically, this object is much closer to a
\emph{stateful constitutive point model} than to a material instance in the
physical sense.  In other words, it already lives at the Gauss-point level,
not at the material-definition level.

\subsection{Inelastic relations already contain the integration algorithm}

The intended design suggests a clean split:
\[
\text{constitutive law} \quad + \quad \text{integration strategy}.
\]

However, for nonlinear materials this separation is only nominal at present.
The current \texttt{InelasticUpdate} strategy simply delegates back to the
material relation, while the actual algorithm lives inside:
\begin{itemize}
  \item \texttt{PlasticityRelation};
  \item \texttt{KentParkConcrete};
  \item \texttt{MenegottoPintoSteel};
  \item \texttt{FiberSection}.
\end{itemize}

Therefore the current strategy layer is semantically strong for elastic
materials, but weak for inelastic ones.  The empty
\texttt{FullImplicitBackwardEuler.hh} placeholder makes this explicit.

\subsection{The constitutive relation and its committed state are fused}

For inelastic models, the relation object itself stores internal variables.
This means that:
\begin{itemize}
  \item the constitutive law is not purely a law;
  \item the law is not shareable as an immutable flyweight;
  \item the same class simultaneously represents parameters, algorithmic cache,
        and committed history.
\end{itemize}

This fusion is convenient for a first implementation, but it becomes a real
limitation for:
\begin{itemize}
  \item thread safety;
  \item state-block vectorisation;
  \item fine-grained checkpointing;
  \item large-scale nonlinear dynamics;
  \item constitutive substepping and rollback.
\end{itemize}

\subsection{The kinematic state is duplicated and over-recorded}

\texttt{MaterialState} stores the current kinematic measure, and for inelastic
materials \texttt{MemoryState} stores the \emph{entire history} of that measure.
At the same time, the inelastic relations store their own internal variables.

This produces two problems:
\begin{itemize}
  \item semantic duplication: committed constitutive history already exists in
        the relation state;
  \item memory amplification: the default inelastic path records every past
        strain state whether or not the algorithm requires it.
\end{itemize}

For nonlinear dynamics or long static load histories this is not sustainable.
The default site state should be:
\begin{itemize}
  \item current trial state,
  \item committed state,
  \item optionally converged last-step state,
\end{itemize}
not an ever-growing vector of all past strains.

\subsection{The flyweight story is currently contradictory}

The comments in \texttt{MaterialInstance} describe the relation as a shared
flyweight.  However, the copy constructor performs a deep copy of the relation.
Because elements and Gauss-point carriers are usually created by copying a
\texttt{Material<Policy>}, the practical behaviour is:
\begin{itemize}
  \item elastic relations are often duplicated per point;
  \item inelastic relations are also duplicated per point;
  \item the claimed flyweight benefit is therefore much weaker than it first
        appears.
\end{itemize}

This is not just a micro-optimisation issue.  It means the current naming and
commentary no longer match the real ownership semantics.

\subsection{Section laws are currently mixed with material laws}

\texttt{TimoshenkoBeamSection3D}, \texttt{MindlinShellSection}, and
\texttt{FiberSection3D} live inside the materials module and satisfy the same
top-level concepts as bulk material laws.  That is convenient, but it masks an
important distinction:
\begin{itemize}
  \item a bulk constitutive law maps local continuum measures to continuum
        stresses;
  \item a section constitutive law maps reduced section measures to section
        resultants;
  \item a fiber section is in fact a \emph{reduction policy plus numerical
        integration over constituents}.
\end{itemize}

This distinction matters if the same section reconstruction is later reused for
scale transfer to a continuum model.

\subsection{Post-processing adapters are embedded in the core wrapper}

\texttt{Material<Policy>} currently knows how to export:
\begin{itemize}
  \item \texttt{InternalFieldSnapshot};
  \item \texttt{SectionConstitutiveSnapshot};
  \item fiber field snapshots.
\end{itemize}

This makes the wrapper convenient for VTK export, but it also means the
central constitutive wrapper now depends on post-processing concerns.  The
core constitutive path should not grow indefinitely through
\texttt{if constexpr}-based export detection.

\subsection{Finite-strain constitutive models live outside the material module}

Hyperelastic relations currently live in \texttt{src/continuum}.  This is
architecturally understandable, because finite-strain kinematics and tensor
operations are housed there, but it also fragments the constitutive story:
\begin{itemize}
  \item small-strain and section laws are described in \texttt{materials};
  \item finite-strain laws are described in \texttt{continuum}.
\end{itemize}

This split becomes problematic when the target is large-displacement,
inelastic, reinforced-concrete analysis.  At that point the constitutive
module and the kinematic module must meet through a cleaner interface.

\section{Performance Implications}

\subsection{Per-fiber type erasure is expensive}

The current \texttt{Fiber} type stores:
\[
\texttt{Material<UniaxialMaterial>}
\]
per fiber.  This gives maximal heterogeneity and a pleasant API, but it is a
poor default for large reinforced-concrete simulations because:
\begin{itemize}
  \item each fiber carries a heap-owning type-erased wrapper;
  \item copies of fiber sections duplicate that structure repeatedly;
  \item the hottest loop in \texttt{FiberSection::integrate\_fibers()} pays
        dynamic-dispatch and scattered-memory costs.
\end{itemize}

For a few dozen fibers this is fine.  For layered shells, detailed RC walls,
or multiscale transfer problems it will quickly become one of the dominant
costs.

\subsection{The default inelastic storage is not HPC-friendly}

\texttt{MemoryState<std::vector<S>>} as the default carrier for inelastic
materials is not suitable as the long-term production layout.  It implies:
\begin{itemize}
  \item heap growth per point;
  \item no fixed-size state block;
  \item poor cache locality;
  \item no easy vectorisation across sites.
\end{itemize}

For HPC-oriented nonlinear analysis the natural constitutive storage unit is a
small fixed record:
\[
\text{trial state} \oplus \text{committed state} \oplus \text{internal variables}.
\]

\subsection{The hot path should remain statically typed}

The present design already suggests the right rule:
\begin{center}
\emph{runtime polymorphism should live at heterogeneous container boundaries,}\\
\emph{not inside the innermost constitutive loops.}
\end{center}

Therefore:
\begin{itemize}
  \item element-level heterogeneous containers may continue to use type
        erasure;
  \item Gauss-point kernels, return mapping, and fiber loops should preferably
        operate on fully typed laws and fully typed state blocks.
\end{itemize}

\section{A Cleaner Separation: Material vs Constitutive Relation}

\subsection{Recommended ontology}

The recommended semantic split is:
\begin{enumerate}
  \item \textbf{Material data}: immutable physical parameters of a phase,
        mixture, or section family.
  \item \textbf{Constitutive law}: the mathematical law that defines stress,
        tangent, and evolution equations for a given measure space.
  \item \textbf{Integration algorithm}: the numerical procedure that advances
        the law in time or load.
  \item \textbf{Constitutive state}: committed and trial internal variables at
        one site.
  \item \textbf{Constitutive site}: the object bound to an integration point
        or section point, owning its state.
  \item \textbf{Reduction policy}: the map from lower-scale fields to reduced
        section measures, and conversely the reconstruction back to a richer
        field when needed.
\end{enumerate}

\subsection{What should be called ``material''}

In this ontology, \emph{material} should denote the physical object or
parameter block, not the entire runtime wrapper.  Examples:
\begin{itemize}
  \item \texttt{ConcreteMaterialData};
  \item \texttt{SteelMaterialData};
  \item \texttt{IsotropicElasticData};
  \item \texttt{RCSectionData}.
\end{itemize}

\subsection{What should be called ``constitutive law''}

A constitutive law should be a purely local mathematical object:
\[
(\text{state},\ \text{kinematics},\ \text{context})
\;\mapsto\;
(\text{response},\ \text{algorithmic tangent},\ \text{new state}).
\]

Its contract should not imply geometry, meshing, or ownership of an
integration point.

\subsection{What should be called ``reduction policy''}

Structural section models are not materials in the same sense as steel or
concrete.  They are reductions of a lower-scale description.  Therefore
\texttt{FiberSection3D} is better understood as:
\begin{center}
\emph{a beam-section reduction policy together with a constitutive quadrature
over constituent fibers.}
\end{center}

The same idea generalises naturally to:
\begin{itemize}
  \item layered shells;
  \item smeared reinforcement layers;
  \item shell-through-thickness integration of concrete damage/plasticity;
  \item scale transfer from beam/shell models to continuum initial fields.
\end{itemize}

\section{A More Scalable Constitutive Contract}

\subsection{Constitutive context}

The current strategy interface only sees the kinematic measure \(k\).  For the
next generation of constitutive models, the local update will need a richer
context:
\[
\mathcal{C}_{n+1}
=
\{\Delta t,\ t_{n+1},\ \text{load increment},\ \text{temperature},
\ \text{kinematic frame},\ \ell_\text{int},\ \text{regularization data}\}.
\]

This is essential for:
\begin{itemize}
  \item rate dependence and viscoplasticity;
  \item nonlinear dynamics;
  \item finite-strain constitutive updates;
  \item nonlocal damage or gradient regularisation;
  \item fracture-energy based softening.
\end{itemize}

\subsection{Suggested compile-time split}

An architecture compatible with the current design philosophy would be:
\begin{align*}
\texttt{ConstitutiveSpace}
&\quad \text{measure types only},\\
\texttt{ConstitutiveLaw}
&\quad \text{physics only},\\
\texttt{IntegrationAlgorithm}
&\quad \text{numerical update only},\\
\texttt{ConstitutiveState}
&\quad \text{small fixed state blocks},\\
\texttt{ConstitutiveSite<Law,Algorithm,Storage>}
&\quad \text{site-level wrapper}.
\end{align*}

This preserves static polymorphism because all of these remain templates or
concept-constrained value types.

\section{Implications for Reinforced Concrete Across Scales}

\subsection{Uniaxial constituent scale}

At the constituent scale the current library already has a credible basis:
\begin{itemize}
  \item \texttt{KentParkConcrete} for confined and unconfined concrete;
  \item \texttt{MenegottoPintoSteel} for reinforcing steel;
  \item \texttt{FiberSection3D} for reduced beam sections.
\end{itemize}

This scale is the correct place for:
\begin{itemize}
  \item section-level hysteresis;
  \item bar yielding and cyclic effects;
  \item confined versus cover concrete distinction;
  \item section-force reconstruction.
\end{itemize}

\subsection{Layered shells}

The next natural scale extension is not yet 3D fracture, but layered shell
constitutive models:
\begin{itemize}
  \item concrete layers through thickness;
  \item membrane reinforcement layers;
  \item optional shear layers;
  \item section resultants obtained by through-thickness integration.
\end{itemize}

Architecturally this should reuse the same ideas as \texttt{FiberSection3D},
but on a shell reduction space rather than a beam section space.

\subsection{Continuum RC with cracking and fracture}

At the continuum scale, merely adding local softening to a stress-strain law is
not enough.  Once cracking localises, the formulation becomes mesh-sensitive
unless an objective regularisation is introduced.  The material architecture
must therefore distinguish:
\begin{itemize}
  \item \textbf{local constitutive law}: plasticity, damage, plastic-damage,
        orthotropic cracking, cohesive traction-separation;
  \item \textbf{regularisation policy}: crack-band scaling, nonlocal averaging,
        gradient enhancement, phase-field, or interface insertion.
\end{itemize}

This distinction is crucial.  Fracture is not only a constitutive issue; it is
also a localisation-regularisation issue.

\subsection{Large displacements and dynamics}

For large-displacement RC, the constitutive layer should not be tied to one
specific small-strain measure.  Instead, the local law should operate on the
appropriate conjugate space supplied by the kinematic formulation:
\begin{itemize}
  \item infinitesimal strain / Cauchy-like stress for small-strain models;
  \item Green--Lagrange / second Piola--Kirchhoff for total Lagrangian models;
  \item current-frame measures for updated-Lagrangian or corotational models.
\end{itemize}

The constitutive law and the kinematic policy must meet through an explicit
measure-space contract, not through ad hoc branches.

\section{Suggested Refactor Path}

\subsection{Stage 1: semantic cleanup}

The first stage should be almost entirely architectural and low-risk:
\begin{itemize}
  \item rename \texttt{MaterialPolicy} toward
        \texttt{ConstitutiveSpacePolicy} or similar;
  \item document that \texttt{Material<Policy>} is an operational wrapper, not
        a physical material;
  \item split bulk laws, section laws, and reduction policies into clearly
        named submodules.
\end{itemize}

\subsection{Stage 2: state redesign}

The second stage should replace \texttt{MemoryState} as the default inelastic
carrier.  The target storage is:
\begin{itemize}
  \item fixed-size committed state;
  \item fixed-size trial state;
  \item optional recorder or history logger outside the hot path.
\end{itemize}

\subsection{Stage 3: genuine algorithm-law separation}

The third stage should make the integration strategy real for nonlinear laws:
\begin{itemize}
  \item \texttt{BackwardEuler};
  \item substepping decorators;
  \item line-search or local damping wrappers;
  \item possibly explicit and operator-split strategies for specific models.
\end{itemize}

At that point the nonlinear relation becomes a law again, not a law plus its
only admissible update algorithm.

\subsection{Stage 4: typed hot path, erased boundary}

The fourth stage should move the runtime-erased wrapper outward:
\begin{itemize}
  \item typed laws and typed states inside Gauss-point loops;
  \item erased wrappers only where heterogeneous containers are genuinely
        required.
\end{itemize}

For fiber sections, the natural next optimisation is:
\begin{itemize}
  \item homogeneous fiber blocks for concrete patches and bar groups;
  \item structure-of-arrays storage for fiber coordinates, areas, and state
        blocks;
  \item optional erased fallback only for rare heterogeneous mixtures.
\end{itemize}

\subsection{Stage 5: RC multiscale and fracture}

Only after the semantic and state cleanup is stable does it make sense to add:
\begin{itemize}
  \item layered RC shell sections;
  \item continuum plastic-damage or smeared cracking models;
  \item fracture-energy regularisation;
  \item cohesive interfaces or phase-field fracture.
\end{itemize}

Without the earlier cleanup, those features would be forced into already
overloaded types and the module would become harder to maintain, not easier.

\section{Non-owning Type Erasure in \texttt{Material.hh}}

\subsection{Viability assessment}

The non-owning type-erasure pattern is viable in the materials module, but not
as a wholesale replacement of the current owning wrapper.  It is viable
precisely because the borrowed erased model stores only:
\begin{itemize}
  \item a pointer to the concrete constitutive object;
  \item a pointer to the integration strategy;
  \item one virtual-table pointer.
\end{itemize}

Therefore the non-owning erased model has a stable footprint of three machine
pointers and fits naturally into a fixed local buffer.  This is fundamentally
different from the owning case, where the concrete constitutive objects are
often too large and too heterogeneous for small-buffer optimisation to pay
off.

\subsection{What problem it solves}

The borrowed view solves a specific architectural problem:
\begin{center}
\emph{how to pass constitutive objects through generic code without cloning or
re-owning them.}
\end{center}

That is useful for:
\begin{itemize}
  \item algorithms that only need a temporary constitutive handle;
  \item post-processing and reconstruction utilities;
  \item adapters layered on top of \texttt{MaterialPoint} and
        \texttt{MaterialSection};
  \item future heterogeneous orchestration code that should not trigger deep
        copies accidentally.
\end{itemize}

\subsection{What problem it does not solve}

The borrowed view does \emph{not} solve constitutive storage by itself.  In
particular, it should not be misused to pretend that the following problem has
disappeared:
\begin{center}
\emph{each constitutive site still needs well-defined ownership of its state.}
\end{center}

This matters especially for reinforced-concrete fiber sections.  Each fiber is
not merely a material label; it is a constitutive site with its own history.
Replacing per-fiber owning storage with a trivial shared borrowed view would
collapse independent histories and would therefore be semantically wrong.

\subsection{Why the strategy belongs in \texttt{Material.hh}}

The strategy object is the correct injection point for the constitutive
integration algorithm because it sits exactly at the operational boundary
between:
\begin{itemize}
  \item the constitutive relation as a law;
  \item the constitutive site as a stateful runtime object.
\end{itemize}

This is the right location to host future algorithms such as:
\begin{itemize}
  \item backward-Euler return mapping;
  \item substepping;
  \item line-search or local damping decorators;
  \item explicit or operator-split updates for dynamics.
\end{itemize}

By contrast, embedding all integration logic directly inside the relation class
fuses law and algorithm and makes it much harder to reuse the same physical law
with different numerical update procedures.

\subsection{Refactor adopted in the current code}

The present refactor therefore adopts a hybrid design:
\begin{itemize}
  \item \texttt{Material<Policy>} remains the owning, heap-backed erased
        wrapper;
  \item \texttt{MaterialConstRef<Policy>} is a borrowed read-only erased view;
  \item \texttt{MaterialRef<Policy>} is a borrowed mutable erased view.
\end{itemize}

This preserves the current ownership semantics of \texttt{MaterialPoint},
\texttt{MaterialSection}, and \texttt{FiberSection}, while enabling generic
algorithms to work with borrowed constitutive handles when ownership transfer
is not intended.

\subsection{Architectural consequence}

The key consequence is that the module now has a clearer boundary:
\begin{itemize}
  \item ownership and lifetime remain explicit;
  \item the hot constitutive storage is not forced into an unsuitable
        small-buffer abstraction;
  \item the library gains a lightweight non-owning polymorphic layer exactly
        where it is semantically valid.
\end{itemize}

\section{Final Opinion}

The current materials module is already much better than a monolithic
``stress-strain class'' design.  It has:
\begin{itemize}
  \item strong measure-based abstractions;
  \item useful concept constraints;
  \item a credible path for structural nonlinear RC through fibers;
  \item a polymorphic envelope that works across solids, beams, and shells.
\end{itemize}

But the module is now at the point where its first successful generalisation is
becoming its main semantic limitation.  The primary issue is no longer whether
the constitutive code works.  It does.  The primary issue is that the present
API still uses the word ``material'' for too many different ontological
objects.

The recommended direction is therefore not a rewrite, but a disciplined
separation:
\begin{center}
\textbf{material data}
\(\rightarrow\)
\textbf{constitutive law}
\(\rightarrow\)
\textbf{integration algorithm}
\(\rightarrow\)
\textbf{stateful constitutive site}
\(\rightarrow\)
\textbf{reduction / reconstruction policy}.
\end{center}

If this separation is carried through while preserving the existing
compile-time contracts, \texttt{fall\_n} will have a robust basis for:
\begin{itemize}
  \item nonlinear reinforced concrete frame analysis;
  \item layered and reduced-order structural models;
  \item finite-strain and dynamic constitutive updates;
  \item objective continuum cracking and fracture formulations.
\end{itemize}

That is the right architectural target.
